\documentclass[10pt,a4paper]{amsart}
\usepackage[margin=19mm]{geometry}
\usepackage{amsmath,amssymb,mathtools}
\usepackage[scr=rsfs,cal=euler]{mathalfa}
\usepackage{xfrac}
\usepackage[unicode,hidelinks,bookmarks=false]{hyperref}
\newcommand{\ie}{i.e.\ }
\newcommand{\cf}{cf.\ }
\newcommand{\resp}{respectively\ }
\newcommand{\Subsection}{\subsection}
\providecommand{\nobreakdash}{\nobreak\hbox{-}}
\setlength{\emergencystretch}{2em}
\multlinegap0pt
\usepackage{bm}
\usepackage{bbm}
\usepackage{dsfont}
\usepackage{xspace}
\usepackage{cancel}
\usepackage{mathtools}
\usepackage{amsthm}
\makeatletter
\@ifundefined{lem}{\newtheorem{lem}{Lemma}}{}
\makeatother
\title[Formation of singularities for the relativistic membrane equ]{Formation of singularities for the relativistic membrane equation with radial symmetry}
\author{Lv Cai, Jianli Liu}
\date{Source: arXiv:2508.07895v1 (CC BY 4.0)}
\begin{document}
\maketitle

\noindent\emph{Source and licence.} Formation of singularities for the relativistic membrane equation with radial symmetry. Version arXiv:2508.07895v1 (CC BY 4.0), distributed under the Creative Commons Attribution 4.0 International licence, as recorded in the arXiv licence metadata for this exact version and shown on the abstract page https://arxiv.org/abs/2508.07895v1. Full rights notice: licenses/arxiv-2508.07895-CC-BY-4.0.txt.\par\smallskip

\noindent\textit{Editorial scope.} Counted unit: the Lem whose source label is \texttt{lem u c prop} in the article (source0.tex lines 789--797, source label \texttt{lem u c prop}), delivered together with the article's own complete original proof of it (source0.tex lines 800--933, one \texttt{proof} environment of 134 lines). No mathematics is written, repaired or re-derived: statement and proof are the article's own text line for line. Required context retained uncounted: def of F (lines 262--266); u v eq3 (lines 361--372); ID (lines 385--389); u c pm (lines 450--461); partial partial v (lines 603--614); coefficient a (lines 616--624); eq of partial R (lines 632--643); coefficient tilde a (lines 645--652); fig1 (lines 722--772). Every definition, display and label of those lines is retained unchanged. Omitted from this package and not redistributed: 1--261, 267--360, 373--384, 390--449, 462--602, 615--615, 625--631, 644--644, 653--721, 773--788, 798--799, 934--1256. Nothing inside a retained excerpt is omitted or altered: every retained body is delivered exactly as the article's lines are written, so no omission receipt of any kind accompanies this package. Citations: the retained excerpts carry no citation command at all, so no bibliography entry of the article is transcribed and no reference is redistributed. Reference adaptation: none was needed and none was made; every reference inside the delivered unit resolves inside it, so no unresolved reference or citation is produced. Printed slips kept literal: no printed word, symbol, hypothesis or number of the counted statement or of its proof was altered. Layout and class: the article's own class is \texttt{amsart} with packages including graphicx, indentfirst, csquotes, amssymb, amsthm, amsmath, xcolor, paralist, hyperref, titlesec, fancyhdr, etoolbox, lipsum; the wrapper is the lane's pinned \texttt{amsart} editorial wrapper at 10pt on A4, which restates the theorem-like environments the retained text needs and adds the article's own macro definitions that the retained text needs (none), with extra wrapper packages (bm, bbm, dsfont, xspace, cancel, mathtools). The delivered numbering is the unit's own. Assets: no retained span contains a figure environment, a graphics-inclusion command, a PDF-page inclusion, a table or a TikZ picture; no image file is copied into this package, the package declares no asset dependency and no image file is redistributed with it. Licence: version arXiv:2508.07895v1 of this article is distributed under the Creative Commons Attribution 4.0 International licence, as recorded in the arXiv licence metadata for this exact version and shown on the abstract page https://arxiv.org/abs/2508.07895v1. A full rights notice is delivered at licenses/arxiv-2508.07895-CC-BY-4.0.txt. Characters: every retained excerpt is pure ASCII, so the delivered engine, XeLaTeX with its default Latin Modern text font, reports no missing character.
	\begin{equation} \label{def of F}
		\begin{aligned}
			\mathcal{F}(u,v) \triangleq  \frac{1}{2}[ ( v^{2}-u^{2}v^{2}-1 ) - \sqrt{(v^{2}-u^{2}v^{2}-1)^2 -4u^{2}v^{2} }  ].
		\end{aligned}
	\end{equation}

\begin{equation} \label{u v eq3}
	\begin{aligned}
		\begin{split}
			\left\lbrace
			\begin{array}{lr}
				u_{t} + uu_{r} + v^{-3}v_{r} = - r^{-1} v^{-2} \mathcal{F} (u,v) & \\ 
				v_{t} + v u_{r} + uv_{r} =- r^{-1} uv . &
			\end{array}	
			\right.
		\end{split}	
	\end{aligned}
\end{equation}

\begin{equation} \label{ID}
	\begin{aligned}
		(u,v) (r,0) = (\bar{u},v_0)(r), \ \ r \in [r_1,r_2].
	\end{aligned}
\end{equation}

\begin{equation} \label{u c pm}
	\begin{aligned}
		\begin{split}
			\left\lbrace
			\begin{array}{lr}	
				\partial_{+} u = -v^{-2}\partial_{+} v - r^{-1}( uv^{-1}+ v^{-2}\mathcal{F}(u,v)) \\
				\partial_{-} u = v^{-2}\partial_{-} v + r^{-1}( uv^{-1}- v^{-2}\mathcal{F}(u,v)).
			\end{array}	
			\right.
		\end{split}	
	\end{aligned}
\end{equation}

\begin{equation} \label{partial partial v}
	\begin{aligned}
		\begin{split}
			\left\lbrace
			\begin{array}{lr}	
				\partial_{+} \partial_{-}v =2 v^{-1} \partial_{+}v \cdot \partial_{-} v + \frac{a_{11}}{2r} \partial_{-}v+ \frac{a_{12}}{2r}\partial_{+}v + \frac{a_{13}}{r^2}, \\
				\partial_{-} \partial_{+}v = 2 v^{-1} \partial_{+}v \cdot \partial_{-} v + \frac{a_{21}}{2r}\partial_{+}v+ \frac{a_{22}}{2r}\partial_{-}v + \frac{a_{23}}{r^2},
			\end{array}	
			\right.
		\end{split}	
	\end{aligned}
\end{equation}

\begin{equation} \label{coefficient a} 
	\begin{aligned}
		a_{11} & =3u-v^{-1}+2v^{-1} \mathcal{F}- v^{-2} \partial_{u} \mathcal{F} - \partial_{v} \mathcal{F}, \\
		a_{12} & =u+v^{-1}-2v^{-1} \mathcal{F}- v^{-2} \partial_{u} \mathcal{F} + \partial_{v} \mathcal{F}, \\
		a_{21} & =3u+v^{-1}-2v^{-1} \mathcal{F}- v^{-2} \partial_{u} \mathcal{F} + \partial_{v} \mathcal{F}, \\
		a_{22} & =u - v^{-1} + 2v^{-1} \mathcal{F}- v^{-2} \partial_{u} \mathcal{F} - \partial_{v} \mathcal{F}, \\
		a_{13}&  = a_{23} = uv(2u-v^{-2}\partial_{u} \mathcal{F}).
	\end{aligned}
\end{equation}	

\begin{equation} \label{eq of partial R}
	\begin{aligned}
		\begin{split}
			\left\lbrace
			\begin{array}{lr}	
				\partial_{+} \tilde{R}_{-} = 2 v^{-1}\tilde{R}_{+}\tilde{R}_{-} + \frac{\tilde{a}_{11}}{2r} \tilde{R}_{-}+ \frac{\tilde{a}_{12}}{2r}\tilde{R}_{+} + \frac{\tilde{a}_{13}}{r^2}, \\
				\partial_{-} \tilde{R}_{+} = 2 v^{-1}\tilde{R}_{+}\tilde{R}_{-} + \frac{\tilde{a}_{21}}{2r}\tilde{R}_{+}+ \frac{\tilde{a}_{22}}{2r}\tilde{R}_{-} + \frac{\tilde{a}_{23}}{r^2}, 
			\end{array}	
			\right.
		\end{split}	
	\end{aligned}
\end{equation}

\begin{equation} \label{coefficient tilde a} 
	\begin{aligned}
		\tilde{a}_{12} & = 2\alpha ( 2v^{-1} \mathcal{F} + v^{-2}  \partial_{u} \mathcal{F} - \partial_{v} \mathcal{F})  + {a}_{12}, \\
		\tilde{a}_{22} & = 2\alpha (2v^{-1} \mathcal{F} - v^{-2}  \partial_{u} \mathcal{F} - \partial_{v} \mathcal{F}) + {a}_{22},\\
		\tilde{a}_{13} & = \alpha^2 \mathcal{F} \cdot ( 2v^{-1} \mathcal{F} + v^{-2}  \partial_{u} \mathcal{F} - \partial_{v} \mathcal{F} )  + \alpha [(3u + 2v^{-1}) \mathcal{F} + uv^{-1} \partial_{u} \mathcal{F} ] +a_{13}, \\
		\tilde{a}_{23} & =\alpha^2 \mathcal{F} \cdot ( 2v^{-1} \mathcal{F} - v^{-2}  \partial_{u} \mathcal{F} - \partial_{v} \mathcal{F} )  + \alpha [(3u -v^{-1}) \mathcal{F} - uv^{-1} \partial_{u} \mathcal{F} ] +a_{23}.
	\end{aligned}
\end{equation}	

\begin{picture}(102.86,67.3)(-15,0)
	\label{fig1}
	\put(8.11,19.05){\vector(1,0){85}}
	\put(12.61,16.05){\vector(0,1){35}}
	%\dottedline(13.11,40.05)(89.61,40.05)
	\multiput(13.04,39.98)(.99351,0){78}{{\rule{.4pt}{.4pt}}}
	%\end
	\put(53.36,36.3){\line(0,1){0}}
	\put(10.08,50){\makebox(0,0)[cc]{$t$}}
	\put(10.08,16.97){\makebox(0,0)[cc]{$0$}}
	\put(89,16.44){\makebox(0,0)[cc]{$r$}}
	\put(20.86,16.44){\makebox(0,0)[cc]{$r_1$}}
	\qbezier(20.51,19.09)(34.82,37.3)(44.19,39.95)
	\qbezier(53.74,18.91)(75.75,37.74)(86.8,39.95)
	\qbezier(33.06,19.09)(38.36,27.93)(45.78,34.65)
	\qbezier(45.96,34.65)(43.22,29.79)(40.84,18.91)
	\put(10.25,40.13){\makebox(0,0)[cc]{$T$}}
	\put(35.36,28.28){\makebox(0,0)[cc]{$C_+$}}
	\put(25.46,32.7){\makebox(0,0)[cc]{$C_+$}}
	%\dottedline(36.77,19.09)(36.95,19.09)
	\multiput(36.7,19.02)(.088,0){3}{{\rule{.4pt}{.4pt}}}
	%\end
	%\dottedline(36.95,19.09)(36.59,18.91)
	\multiput(36.88,19.02)(-.177,-.088){3}{{\rule{.4pt}{.4pt}}}
	%\end
	%\dashline{1}(45.61,34.47)(36.77,19.09)
	\multiput(45.54,34.4)(-.033228,-.057817){14}{\line(0,-1){.057817}}
	\multiput(44.61,32.78)(-.033228,-.057817){14}{\line(0,-1){.057817}}
	\multiput(43.68,31.16)(-.033228,-.057817){14}{\line(0,-1){.057817}}
	\multiput(42.75,29.54)(-.033228,-.057817){14}{\line(0,-1){.057817}}
	\multiput(41.82,27.93)(-.033228,-.057817){14}{\line(0,-1){.057817}}
	\multiput(40.89,26.31)(-.033228,-.057817){14}{\line(0,-1){.057817}}
	\multiput(39.96,24.69)(-.033228,-.057817){14}{\line(0,-1){.057817}}
	\multiput(39.02,23.07)(-.033228,-.057817){14}{\line(0,-1){.057817}}
	\multiput(38.09,21.45)(-.033228,-.057817){14}{\line(0,-1){.057817}}
	\multiput(37.16,19.83)(-.033228,-.057817){14}{\line(0,-1){.057817}}
	%\end
	%\dashline{1}(36.77,19.09)(36.77,19.09)
	\put(36.7,19.02){\line(0,1){0}}
	%\end
	\put(47,36.42){\makebox(0,0)[cc]{$P$}}
	\put(47,26.16){\makebox(0,0)[cc]{$C_-$}}
	\put(38.5,21.74){\makebox(0,0)[cc]{$\Omega_P$}}
	\put(30.5,16.44){\makebox(0,0)[cc]{$P_+$}}
	\put(36.59,16.44){\makebox(0,0)[cc]{$P_0$}}
	\put(43,16.44){\makebox(0,0)[cc]{$P_-$}}
	\put(55.86,32){\makebox(0,0)[cc]{$\Lambda(T)$}}
	\put(72.83,27.4){\makebox(0,0)[cc]{$C_-$}}
	\put(53.21,16.44){\makebox(0,0)[cc]{$r_2$}}
	\put(21,7){Figure 3.2: Domains $\Lambda(T)$ and $\Omega_P$}
\end{picture}

\begin{lem} \label{lem u c prop}
	Suppose that system \eqref{u v eq3} with initial data \eqref{ID} admits a classical solution in $ \Lambda (T)$ for some $T>0$. Then under the assumptions (A1)-(A2), the solution satisfies 
	\begin{equation} \label{lem R}
		\begin{aligned}
			u_2 \leq u \leq u_1, \ \tilde{R}_{\pm} > 0 ,
		\end{aligned}
	\end{equation}	
	in $\Lambda (T)$, where $	\tilde{R}_{\pm} = \partial_{\pm} v - \frac{\alpha\mathcal{F}}{r}$ for $\alpha \geq 1$.
\end{lem}

\begin{proof}
	We shall prove this lemma by the argument of continuity. First, we will check the estimate \eqref{lem R} is valid for the initial data. Noting the second equation of \eqref{u v eq3}, one has
	\begin{equation*} 
		\begin{aligned}
			v_{t} = -uv_{r} - vu_r-\frac{uv}{r}.
		\end{aligned}
	\end{equation*}	
	Hence,
	\begin{equation*} 
		\begin{aligned}
			\partial_{\pm} v  = v_{t} + (u \pm v^{-1}) v_{r}  = \pm v^{-1} v_{r} - vu_r-\frac{uv}{r}.
		\end{aligned}
	\end{equation*}	
	According to (A2), we have
	\begin{equation*} 
		\begin{aligned}
			\tilde{R}_{\pm}(r,0) =  -v_0 \left(   \bar{u}^{\prime}(r) + \frac{\bar{u} + \alpha \mathcal{F}_{0}}{r} \right) >0 ,
		\end{aligned}
	\end{equation*}	
	Therefore, the estimates of \eqref{lem R} hold on $\{(r,t)|t=0, r_1 \leq r \leq r_2 \}$.
	
	
	Let $P$ be an arbitrary point in the area $\Lambda (T)$. The backward $C_{+}$ and $C_{-}$ characteristic curves are drawn from the point $P$ intersecting the $r$-axis with the points $P_{+}$ and $P_{-}$, respectively. The backward $C_{0}$ curve passing through $P$ intersects the $r$-axis at a point $P_{0}$. We denote by $\Omega_{P}\subset \Lambda (T)$ a closed triangle domain
	bounded by $\widehat{P_{+} P}$, $\widehat{P_{-} P}$	and $\widehat{P_{+} P_{-}}$, see Fig. \ref{fig1}. In the following we shall prove that if the inequalities in \eqref{lem R} hold
	for every point in $\Omega_{P} \backslash \{P\}$, then they also hold at $P$. The proof of this assertion will be given in three steps.
	
	{\textit{Step 1}}.	By the assumption that the estimates in \eqref{lem R} hold for any point in $\Omega_{P} \backslash \{P\}$, we have $\partial_{\pm}v> 0$ in $\Omega_{P} \backslash \{P\}$. Noting \eqref{u c pm}, along the characteristics $C_{+}$ and $C_{-}$, we have
	\begin{equation*} 
		\begin{aligned}
			\partial_{+} u  < 0 \ \ \text{on } \widehat{P_{+} P} \ \ \ \ \text{and} \ \ \ \	\partial_{-} u   > 0 \ \ \text{on } \widehat{P_{-} P}.
		\end{aligned}
	\end{equation*}	
	Thus we obtain 
	\begin{equation*} 
		\begin{aligned}
			u_2 < u(P_{-}) < u(P) < u(P_{+}) < u_1.
		\end{aligned}
	\end{equation*}	
	
	{\textit{Step 2}}. We shall analyze the signitures of coefficients in coupled equations \eqref{partial partial v} and \eqref{eq of partial R} in this step. By \eqref{def of F} and the direct computation, one can get
	\begin{equation} \label{partial u F}
		\begin{aligned}
			\partial_{u}\mathcal{F} & =  \frac{1}{2}\left[ -2uv^2 - \frac{2(v^2-u^{2}v^2-1) ( -2uv^2 ) - 8 uv^2 }{2\sqrt{(v^2-u^{2}v^2-1)^2 -4u^{2}v^{2} }} \right] \\
			& =  -uv^{2} \left(  1- \frac{v^2-u^{2}v^2+1 }{\sqrt{(v^2-u^{2}v^2-1)^2 -4u^{2}v^{2} }} \right) <0 ,
		\end{aligned}
	\end{equation}
	and
	\begin{equation}
		\begin{aligned}
			\partial_{v}\mathcal{F}& =  \frac{1}{2}\left[ 2v(1-u^2) - \frac{2(v^2-u^{2}v^2-1) \cdot(-2v)(1-u^2)  - 8 u^2 v }{2\sqrt{(v^2-u^{2}v^2-1)^2 -4u^{2}v^{2} }} \right]  \\
			& = v \left[ (1-u^2) - \frac{v^2(1-u^{2})^{2} - (1+u^2) }{\sqrt{(v^2-u^{2}v^2-1)^2 -4u^{2}v^{2} }} \right].
		\end{aligned}
	\end{equation}
	Hence, by \eqref{coefficient a} we have
	\begin{equation*} 
		\begin{aligned}
			a_{12} & =u+v^{-1}-2v^{-1}\mathcal{F}- v^{-2} \partial_{u}\mathcal{F} + \partial_{v}\mathcal{F} \\
			& =u+v^{-1}- v^{-1} \left[ (v^2-u^{2}v^2-1)-\sqrt{(v^2-u^{2}v^2-1)^2 -4u^{2}v^{2} }\right]  \\
			& \ \ \ \ + u \left(   1- \frac{v^2-u^{2}v^2+1 }{\sqrt{(v^2-u^{2}v^2-1)^2 -4u^{2}v^{2} }}  \right)\\
			& \ \ \ \   + v \left[ (1-u^2) - \frac{v^2(1-u^{2})^{2} - (1+u^2) }{\sqrt{(v^2-u^{2}v^2-1)^2 -4u^{2}v^{2} }} \right]   \\
			& = (u+v^{-1}) \left( 2- \sqrt{\frac{v^2-u^{2}v^2-1+2uv}{v^2-u^{2}v^2-1-2uv}} \right) ,
		\end{aligned}
	\end{equation*}	
	and
	\begin{equation*} 
		\begin{aligned}
			a_{22} & =u-v^{-1}+2v^{-1}\mathcal{F}- v^{-2} \partial_{u}\mathcal{F} - \partial_{v}\mathcal{F} \\
			& =u-v^{-1}+ v^{-1} \left[ (v^2-u^{2}v^2-1)-\sqrt{(v^2-u^{2}v^2-1)^2 -4u^{2}v^{2} }\right]  \\
			& \ \ \ \ + u \left(   1- \frac{v^2-u^{2}v^2+1 }{\sqrt{(v^2-u^{2}v^2-1)^2 -4u^{2}v^{2} }}  \right) \\
			& \ \ \ \  - v \left[ (1-u^2) - \frac{v^2(1-u^{2})^{2} - (1+u^2) }{\sqrt{(v^2-u^{2}v^2-1)^2 -4u^{2}v^{2} }} \right]   \\
			& = (u-v^{-1}) \left( 2- \sqrt{\frac{v^2-u^{2}v^2-1-2uv}{v^2-u^{2}v^2-1+2uv}} \right),
		\end{aligned}
	\end{equation*}	
	where we use the assumption (A1). Obviously, by \eqref{partial u F}, we also have
	\begin{equation*} 
		\begin{aligned}
			a_{13} = a_{23} = uv(2u-v^{-2}\partial_{u}\mathcal{F})> 0.
		\end{aligned}
	\end{equation*}	
	Noting \eqref{coefficient tilde a}, one has
	\begin{equation} 
		\begin{aligned}
			\tilde{a}_{12} & = 2\alpha (2v^{-1}\mathcal{F}+ v^{-2} \partial_{u}\mathcal{F} - \partial_{v}\mathcal{F}) + {a}_{12} \\
			& = 2\alpha (- a_{12} +u +v^{-1}) + {a}_{12}\\
			& = (u +v^{-1}) \left[  (2 \alpha -1) \sqrt{\frac{v^2-u^{2}v^2-1+2uv}{v^2-u^{2}v^2-1-2uv}} - 2 \alpha+2  \right] > 0,
		\end{aligned}
	\end{equation}	
	for any $ \alpha \geq 1$, and
	\begin{equation}
		\begin{aligned}
			\tilde{a}_{22} & = 2\alpha (2v^{-1}\mathcal{F}- v^{-2} \partial_{u}\mathcal{F} - \partial_{v}\mathcal{F}) + {a}_{22} \\
			& = 2\alpha (a_{22} -(u -v^{-1})) + {a}_{22}\\
			& = (u -v^{-1}) \cdot  2 \alpha \left( -  \sqrt{\frac{v^2-u^{2}v^2-1-2uv}{v^2-u^{2}v^2-1+2uv}} +1 \right) \\
			& \ \ \ \ + (u -v^{-1}) \left(  2 - \sqrt{\frac{v^2-u^{2}v^2-1-2uv}{v^2-u^{2}v^2-1+2uv}} \right) > 0,
		\end{aligned}
	\end{equation}	
	for any $ \alpha \geq 0$. Besides, when $ \alpha =1$, one has
	\begin{equation*} 
		\begin{aligned}
			\tilde{a}_{13} &= ( 3u + v^{-1} + 2v^{-1} \mathcal{F} + v^{-2}  \partial_{u} \mathcal{F} - \partial_{v} \mathcal{F} ) \mathcal{F} + 2u^2 v \\
			& = (- a_{12} +4u +2v^{-1})+ 2u^2 v \\
			& = \left[ 2u + (u + v^{-1}) \sqrt{\frac{v^2-u^{2}v^2-1+2uv}{v^2-u^{2}v^2-1-2uv}} \right] \mathcal{F} + 2u^2 v  > 0,
		\end{aligned}
	\end{equation*}	
	and
	\begin{equation*} 
		\begin{aligned}
			\tilde{a}_{23} &= ( 3u + v^{-1} - v^{-2}  \partial_{u} \mathcal{F} - \partial_{v} \mathcal{F} ) \mathcal{F} + 2uv (u - v^{-2}  \partial_{u} \mathcal{F}) \\
			& = \left[ a_{22} + 2u + v^{-1} \left( 3- v^{2} (1-u^{2}) + \sqrt{\frac{v^2-u^{2}v^2-1-2uv}{v^2-u^{2}v^2-1+2uv}} \right) \right] \mathcal{F} \\
			& \ \ \ \ + 2uv (u - v^{-2}  \partial_{u} \mathcal{F})  > 0.
		\end{aligned}
	\end{equation*}	
	Moreover, it is clear that 
	\begin{equation} \label{tilde a13 a23} 
		\begin{aligned}
			\tilde{a}_{13}>0, \ \tilde{a}_{23}>0,
		\end{aligned}
	\end{equation}	
	when $ \alpha > 1$.
	
	
	{\textit{Step 3}}.	Now we prove $\tilde{R}_{+} (P) > 0$. Suppose $\tilde{R}_{+}(P) = 0$ and $\tilde{R}_{-}(P) \geq 0$. Then by the assumption that the inequalities in \eqref{lem R} hold for every point in $\Omega_{P} \backslash \{P\}$, we have $\partial_{-}\tilde{R}_{+} \leq 0$ at $P$. However, by the first equation of \eqref{eq of partial R} and \eqref{tilde a13 a23} at the point $P$, we derive
	\begin{equation*} 
		\begin{aligned}
			\partial_{-} \tilde{R}_{+} & = \frac{\tilde{a}_{12}}{2r} \tilde{R}_{-} + \frac{\tilde{a}_{13}}{r^2} > 0 ,
		\end{aligned}
	\end{equation*}	
	which leads to a contradiction. Hence, we have
	$\tilde{R}_{+}(P) > 0$. Similarly, one can get $\tilde{R}_{-}(P)> 0$. Therefore we finish the proof of the
	assertion. 
	
	Consequently, we obtain the estimate \eqref{lem R}. We complete the proof of this lemma.
	
\end{proof}

\end{document}
